\section{Discussion and Threats to Validity}
\label{sec:threats}
\paragraph{Meaning of the guarantees.}
Safety differs from eventual progress~\cite{safety}. Our retry theorem fixes
corrected intent, source membership, admission predicates, and channel inputs.
Some stability lemmas follow directly from this invariant, exposing what retries
must preserve. New commits, policy edits, or tag rewrites can require replanning;
permanent failure need not converge.

Shared-counter groups can advance prerelease counters on retry and reselect
successful members. They are outside our per-package version-stability and
decreasing-plan convergence proofs.

The model abstracts artifact contents and combines publication and recording
into one transition. Lost responses and failed recording remain integration
obligations; Git locks do not make registries transactional. Own-reason failure
exceptions, recording, and finalization require refinement beyond the abstract
blocked-set rule; selected test traces do not close that gap.

\paragraph{Internal and construct validity.}
The author developed the system, model, oracle, and tests; agreement can preserve
shared mistakes. The restricted oracle omits semantics, and full-profile
categories are heuristic. The prospective matrix was frozen before execution
but designed after the earlier campaign, not independently replicated.
The finite projection covers stable-channel, one-intent states; post-write
faults are responses, not process crashes. Integrity checks establish record
consistency; concurrency tests cover selected interleavings. Histories, package
rows, episodes, fixture/CI repetitions, commands, and assertions remain
separate units.

\paragraph{External validity.}
Synthetic fixtures establish controlled behavior; historical replay reconstructs
policies, and comparisons pin configurations. No user study measures intent
correctness or review effort. The control-history polyrepo fits the model;
multi-history coordination, peer choreography, and their scale remain
unevaluated. Platform scripts, agent gates, and ecosystems demonstrate
integration scope, not general effectiveness. One contended host cannot
establish service-level promises.

\paragraph{Exploratory outreach context.}
Eight retained September review inventories cover at least 211 distinct
repositories after deduplication, with unequal review depth. One pass spans
21 native ecosystems; a separate search's 6,731 hits across 1,512 repositories
are discovery coverage, not reviewed-project or defect counts.
Subjects include RelayCast (9 GitHub stars) and Paperclip (80,381) at the
recorded snapshot; stars describe visibility, not code size.
Three projects illustrate different boundaries. In RelayCast, a two-clone
fixture simulates successful upload followed by a rejected push: no remote tag
is created, yet the script reports success; no production failure is claimed.
Two historical Astro sequences publish npm packages, fail before marketplace
publication, then finish green on retry while skipping those destinations;
later manual completion was not assessed. An Apache Arrow fixture with a
stubbed workflow client reproduces selection of a stale successful run for a
new invocation. These observations motivate destination-specific completion
and invocation identity in fleet recovery. Healthy checks, resolved incidents,
and later rejected claims remain in the records. This selected review neither
estimates defect prevalence nor validates our protocol on these projects.

\section{Related Work}
\label{sec:related}
\paragraph{Release and build systems.}
Release engineering is a distinct research problem~\cite{adams}.
Monorepo practice~\cite{potvin} and build-system models~\cite{mokhov} address
ordering and reuse; incremental builds admit dynamic dependencies~\cite{pluto}.
Our fixed graph records publication obligations without claiming a new
topological algorithm or incremental rebuilding.
Lerna, Nx, Changesets, release-please, and semantic-release provide different
versioning and publication workflows~\cite{lerna,nx,changesets,releaseplease,
semanticrelease}. Section~\ref{sec:compare} compares operational traces of
specified configurations, not formal reductions of those complete tools.
Spinnaker supports parallel stages and execution histories~\cite{spinnaker};
Nix uses input-derived store paths~\cite{nix}; TOSCA standardizes lifecycle
interfaces~\cite{tosca}; OpenGitOps calls for versioned desired state and
reconciliation~\cite{gitops}. Our ledger tracks completed releases across
independent publication targets without replacing these systems' state models.

Versioning studies motivate explicit policy~\cite{raemaekers,bogart,decan},
without proving compatibility. EDOS studies installability~\cite{edos};
CI effects depend on context~\cite{huangci}.
DORA associates greater AI adoption with higher delivery throughput and
instability~\cite{dora2025}. Our accompanying agent guide directs agents to
inspect deterministic plans, retain tests, invoke CI releases, and reconcile
partial outcomes. The engine implements a model with eight conditional
mathematical guarantees; agent effectiveness and full implementation refinement
remain unproven. DORA motivates this design, not its validation.
Workflow duplication motivates shared configuration~\cite{cardoen};
maintenance gains remain unmeasured. Formal VCS models~\cite{swierstraloh,hpt,mimram,mrdt}
and Git-mining cautions~\cite{birdgit} address concerns beyond our release ledger.

\paragraph{Recovery and formal assurance.}
Buildbox formalizes build, install, ship, and update transitions~\cite{liusmith};
Aeolus studies component states, dependencies, and reachability~\cite{aeolus}.
These precede our narrower publication-recovery question. NEPHELE applies
sagas to distributed lifecycle operations~\cite{nephele}; Bui et al.'s
LLM-assisted CI/CD diagnosis and repair~\cite{bui} measures different outcomes.
Sagas~\cite{sagas} and the limits of transaction abstractions~\cite{gray}
provide the recovery setting. Independently committed and idempotently
processable operations~\cite{helland07,helland12} explain why uncertain
outcomes need explicit treatment. Our contribution specializes admission and
discharge to versioned consumers; it is not the invention of forward recovery.
Anvil~\cite{anvil} verifies reconciliation properties of Kubernetes controller
implementations, whereas our handwritten proofs establish a release model and
our tests sample implementation behavior. This is a significant difference
in assurance. Methodological cautions about fuzz evaluation~\cite{klees}
likewise motivate separating core agreement from unsupported generalization.

\section{Conclusion}
\label{sec:conclusion}
Stage prerequisites enforce artifact availability; consumer-based admission
retains obligations after providers succeed. Under the retry invariant,
separate train and fresh-work windows preserve versions and permit convergence.
All 140 mechanism episodes retained expected plans; the source-window
counterfactual lost work in 84/120 fault episodes. Fleet extensions require
repository-qualified histories and evidence of consumer boundaries.
Multi-history recovery, recording failures, and randomized interleavings
remain targets for evaluation and mechanized refinement~\cite{lamport}.

\section*{Data Availability}
Campaign records retain runtime pins, harnesses, raw observations, verified
archives, offline verifiers, and fixture Git metadata. Outreach inventories,
deduplication, and case provenance are retained locally. These records can
reveal authorship; no anonymized package is available, and public release
remains undecided.

\section*{Acknowledgments}
\noindent\textit{AI assistance.}
Claude Fable provided most assistance; Claude Opus, Gemini, and OpenAI Codex
also contributed. Assistance went beyond editing: implementation and experiment
tooling (Sections~\ref{sec:system} and~\ref{sec:eval}), and manuscript drafting
and revision, including the system description, model exposition, evaluation,
and discussion. Executions produced the measurements; the human author
remains responsible for all claims and results.
